\section{Testing and measurements}
\label{sec:quality}

\subsection{Test strategy}

The project has 409 automated tests, written with JUnit and run by a single command,
\code{./mvnw verify}. They are of four kinds.

\begin{description}[leftmargin=1.2em,style=nextline]
  \item[Unit tests of the model.] Each class is tested through its public interface: the
    recursive array in one to six dimensions, the two boundary conditions, the truth table of
    every operator against its definition in the subject, the size and content of every
    predefined neighborhood in dimensions one to five.
  \item[Tests against the subject.] The figures of the subject are used as oracles. The growth
    automaton must reproduce figure 2 row by row; the Sierpi\'nski rule must match the truth table
    of figure 4 \emph{and} the arithmetic formula of the subject on every cell of 31 generations;
    the worked example of section 2.3 must parse to the expected tree. For the Game of Life, the
    classical patterns are checked: a block is stable, a blinker has period 2, a glider moves by
    one cell diagonally every four generations.
  \item[Cross-validation.] Two independent implementations of one-dimensional automata exist: the
    general engine and the fast simulator of the rule explorer. A test generates random rules,
    translates them from one representation to the other, and requires both to agree on every
    cell. A bug in either, or in the translation, would break it.
  \item[End-to-end tests.] The command-line interface is driven with real arguments in a
    temporary directory, and its output, exit status and generated files are checked: this
    exercises the whole chain from an XML file to a PNG image. The Swing windows have their own
    tests, which build the real windows without showing them and press their buttons; they need
    a display and are skipped on a build server.
\end{description}

Invalid inputs receive as much attention as valid ones: 27 kinds of invalid definition files and
18 kinds of syntax errors are each checked for the message shown to the user. Every file of the
\code{examples} directory is loaded and run by a test, so a broken example cannot be shipped.

\begin{table}[H]
  \centering\small
  \begin{tabular}{@{}lrrrrr@{}}
    \toprule
    Package & Classes & Lines & Tests & Line coverage & Branch coverage \\
    \midrule
    \code{grid} & 16 & 345 & 56 & 97\,\% & 94\,\% \\
    \code{neighborhood} & 5 & 148 & 63 & 99\,\% & 94\,\% \\
    \code{rule} (with \code{expression}, \code{parser}) & 28 & 279 & 93 & 100\,\% & 96\,\% \\
    \code{automaton} & 6 & 59 & 19 & 95\,\% & 95\,\% \\
    \code{config} & 16 & 358 & 95 & 97\,\% & 91\,\% \\
    \code{simulation} & 5 & 200 & 18 & 96\,\% & 83\,\% \\
    \code{render}, \code{ui.text} & 5 & 143 & 9 & 99\,\% & 85\,\% \\
    \code{explore} & 6 & 262 & 26 & 98\,\% & 92\,\% \\
    \code{cli} & 11 & 286 & 25 & 94\,\% & 89\,\% \\
    \midrule
    All of the above & 98 & 2\,080 & 404 & 97\,\% & 92\,\% \\
    \code{ui.swing} & 12 & 460 & 5 & \multicolumn{2}{c}{needs a display} \\
    \bottomrule
  \end{tabular}
  \caption{Tests and coverage by package, as measured by JaCoCo in a headless build. Classes
    include nested classes; lines are executable lines.}
  \label{tab:coverage}
\end{table}

\subsection{Tooling}

\begin{itemize}
  \item \textbf{Build.} Maven, through the wrapper, so that the only prerequisite is a JDK. The
    application itself has no dependency: only the tests need JUnit.
  \item \textbf{Compiler as reviewer.} The code is compiled with all the lints of \code{javac}
    enabled and treated as errors: an unchecked conversion or a missing case does not compile.
  \item \textbf{Continuous integration.} A GitHub Actions workflow builds and tests on Linux,
    macOS and Windows, with two versions of the JDK, and validates the example files with the
    packaged application.
  \item \textbf{Version control.} The history follows the Conventional Commits convention, one
    commit per coherent step of the work.
\end{itemize}

\subsection{Performance}

\Cref{tab:performance} gives the time needed to compute one generation. The cost is proportional
to the number of cells times the number of cells the rule reads; large grids use all the cores.
A million-cell Game of Life runs at about 35 generations per second, far beyond what the display
needs, which confirms that keeping the recursive array and an interpreted rule was affordable.

\begin{table}[H]
  \centering\small
  \begin{tabular}{@{}llrr@{}}
    \toprule
    Automaton & Neighborhood & Cells & Time per generation \\
    \midrule
    Sierpi\'nski, 1D & G2 & 257 & 0.02 ms \\
    Game of Life, $64 \times 96$ & G8* & 6\,144 & 0.8 ms \\
    Game of Life, $512 \times 512$ & G8* & 262\,144 & 7 ms \\
    Game of Life, $1024 \times 1024$ & G8* & 1\,048\,576 & 29 ms \\
    Crystal, $65^3$ & G6* & 274\,625 & 8 ms \\
    Clouds, $40^3$ & G26* & 64\,000 & 13 ms \\
    Hypercube, $25^4$ & G8* (4D) & 390\,625 & 18 ms \\
    \bottomrule
  \end{tabular}
  \caption{Median time per generation on an Apple M1~Pro (10 cores), OpenJDK 27. Grids below
    8\,192 cells are computed on a single thread.}
  \label{tab:performance}
\end{table}
